\begin{textsample}{POS Dim 2 – global\_south\_2024\_12 – Score 19.00 – global\_south\_2024\_12/118644850.txt}  \label{ex:f2_pos_209}
Assad falls thanks to a weak Tehran A defaced portrait of late \textbf{Syrian} president Hafez al-Assad stands above a ransacked government security facility, in Damascus.
Islamist-led \textbf{rebels} declared on December 8, that they have taken the \textbf{Syrian} \textbf{capital} in a lightning offensive, sending President Bashar al-Assad fleeing and ending five decades of Baath rule in Syria. ( Photo : AFP ) Summary Iran and its \textbf{proxy} network have been revealed to be paper tigers.
More defeats will surely follow this one.
This is a Mint Premium article gifted to you.
Subscribe to enjoy similar stories.
The fall of Bashar al-Assad ' s regime in Syria marks the eclipse of Iran ' s self-styled \textbf{axis} of resistance in the Middle East.
While the successful advance of the \textbf{Syrian} Sunni insurgents into Damascus seemed sudden, it didn't happen overnight.
Mr.
Assad ' s flight from the country seals what has been a disastrous few months for Tehran and its \textbf{allies}.
By removing the Assad regime, the insurgents -- known as Hayat Tahrir al-Sham -- have cut off Tehran ' s main \textbf{proxy} \textbf{militia}, Hezbollah, from its supply lines.
This was possible because of Israel ' s mauling of Hezbollah in November, prompting the terror group to accept a cease-fire with the Jewish state.
The decision was in direct contrast with Hezbollah ' s goal when it entered the war on Oct. 8, 2023.
As Hassan Nasrallah, its now-deceased leader, expressed at the time :"The Lebanon front will not stop before the \textbf{aggression} on Gaza stops.
The resistance in Lebanon won't stop supporting and assisting the people of Gaza, the West Bank and the oppressed people in these holy lands."By the time both parties agreed to a cease-fire, Nasrallah was dead.
Hezbollah decided that sticking by its junior \textbf{ally} in Gaza was no longer worth the price.
The dramatic events in Syria rapidly followed.
Hezbollah ' s actions came after Iran ' s decision to avoid any major \textbf{retaliation} for Israel ' s Oct. 26 \textbf{attacks} on \textbf{targets} inside Iran.
This reluctance persisted despite Supreme Leader Ali Khamenei ' s promise of a"crushing response to what they are doing to Iran... and to the resistance front"after the \textbf{attack}.
Iran hesitated in part because of the failure of its Iraqi and \textbf{Yemeni} \textbf{proxies} to penetrate Israeli \textbf{air} defenses with \textbf{drones} and \textbf{missiles}.
Israel continued to make progress against Hamas in Gaza, despite Iran ' s various efforts to open support fronts to relieve the pressure.
Gaza ' s Islamists are now isolated, facing Israel alone.
The apparent eclipse of Iran ' s \textbf{proxy} \textbf{axis}, its helplessness in the wake of Israeli countermeasures, and the loss of Syria are of particular note because the alliance had fared so well for so long.
Iran ' s first success in the field of irregular warfare came with its creation of Hezbollah in 1982.
It piloted Hezbollah to a victorious insurgency against Israel from 1985 to 2000, and then to undisputed political domination of Lebanon after 2008.
After the fall of Saddam Hussein in 2003, Tehran-supported Shiite \textbf{militias} became the dominant political-military force in Iraq.
Outreach to Hamas via Hezbollah from the early 1990s secured Iran a foothold in the Israel-Palestinian conflict, with its enormous symbolic importance in the Arab and Muslim worlds.
The 2010 Arab Spring, the 2012 destruction of the Saleh regime in Yemen and Iran ' s rescue of the Assad regime in Syria that same year all contributed to the expansion of \textbf{Iranian} influence across the \textbf{region}.
On the eve of Oct. 7, 2023, Iran ' s \textbf{proxy} system seemed at the height of its power.
Iran, with its formidable ballistic-missile array and its \textbf{nuclear} ambitions, appeared to have succeeded in turning indirect conflict into a tool of state.
A year on, however, everything looks different.
Iran ' s error was that it neglected a basic rule of irregular warfare : Avoid being drawn into the open on the conventional battlefield.
Doughty proponents of irregular warfare, from George Washington to Ho Chi Minh, all lived by this rule.
Iran and its \textbf{allies} have abandoned it and suffered the consequences.
The Oct. 7 \textbf{attack}, in addition to being a massacre, was a strategically absurd decision by a tiny statelet to \textbf{launch} a conventional offensive against a vastly more powerful neighbor.
Hezbollah ' s choice to bombard Israel in the following days constituted a similar misstep.
Iran ' s decision in April this year to abandon \textbf{proxy} warfare and \textbf{launch} a direct \textbf{attack} on Israel compounded the error.
In each case, Israel ' s response laid bare the profound inferiority of the \textbf{Iranians} and their \textbf{allies} in direct confrontation.
The result : Hamas and Hezbollah are decimated, Gaza is a smoking ruin, Southern Lebanon is a pile of rubble, and Iran is exposed as helpless before Israeli \textbf{air} power.
The \textbf{region} now sees Iran and its misnamed \textbf{axis} of resistance as a paper tiger.
The \textbf{rebels} ' assault in Syria and the stunning collapse of the Assad regime are the first fruits of this new look.
More will doubtless come.
Mr.
Spyer is director of research at the Middle East Forum and director of the Middle East Center for Reporting and Analysis.
He is author of"Days of the Fall : A Reporter ' s Journey in the Syria and Iraq Wars."

% matched lemmas: aggression, air, ally, attack, axis, capital, drone, iranian, launch, militia, missile, nuclear, proxy, rebel, region, retaliation, syrian, target, yemeni
\end{textsample}
